\documentclass[12pt]{article}
% -------------------
% Packages
% -------------------
\usepackage{
	amsmath,			% Math Environments
	amssymb,			% Extended Symbols
	amsthm,			    % Theorem Environments
	cancel,			    % Use Cancels
	enumerate,		    % Enumerate Environments
	graphicx,			% Include Images
	lastpage,			% Reference Lastpage
	multicol,			% Use Multi-columns
	multirow,			% Use Multi-rows
	xcolor,			    % Use Colors
	float,
	geometry,
	quiver
	}

% -------------------
% Header & Footer
% -------------------
\usepackage{fancyhdr}

\newcommand{\assignment}[1]{
	\fancypagestyle{title}{
		%Headers
		\fancyhead[L]{Jake Bridge}
		\fancyhead[C]{MATH 418}
		\fancyhead[R]{HW #1}
		\renewcommand{\headrulewidth}{0.2pt}
		%Footers
		\fancyfoot[L]{}
		\fancyfoot[C]{}
		\fancyfoot[R]{}
		\renewcommand{\footrulewidth}{0.0pt}
	}
	\thispagestyle{title}
}
\fancypagestyle{pages}{
	%Headers
	\fancyhead[L]{}
	\fancyhead[C]{}
	\fancyhead[R]{}
	\renewcommand{\headrulewidth}{0.0pt}
	%Footers
	\fancyfoot[L]{}
	\fancyfoot[C]{}
	\fancyfoot[R]{}
	\renewcommand{\footrulewidth}{0.0pt}
}
\headheight=18pt
\footskip=14pt

\pagestyle{pages}

% ---------------
% Page Layout
% ---------------
\geometry{letterpaper,
bindingoffset=0.2in,
left=1in,
right=1in,
top=1in,
bottom=1in,
footskip=.25in}
% -------------------
% Font
% -------------------
%\usepackage[T1]{fontenc}
%\usepackage{charter}

%\usepackage[T1]{fontenc}
%\usepackage{mathpazo}

%\usepackage[bitstream-charter]{mathdesign}
%\usepackage[T1]{fontenc}


% -------------------
% Commands
% -------------------

% Problem Labels
\newcounter{problem}[section]
\newenvironment{problem}[1][\theproblem]{\refstepcounter{problem}\noindent \textbf{Problem~#1.\quad}}{\vskip\smallskipamount}


\newenvironment{solution}{\noindent \textbf{Solution.\quad}}{\vskip\bigskipamount}


% Special Characters
\newcommand{\N}{\mathbb{N}}
\newcommand{\Z}{\mathbb{Z}}
\newcommand{\Q}{\mathbb{Q}}
\newcommand{\R}{\mathbb{R}}
\newcommand{\C}{\mathbb{C}}

\newcommand{\mc}[1]{\mathcal{#1}}

\newcommand{\ob}[1]{\text{ob}(#1)}
\newcommand{\natt}{\Rightarrow}
\newcommand{\iso}{\cong}
\newcommand{\op}{\text{op}}

\newcommand{\Set}{\textbf{Set}}
\newcommand{\Hom}{\textbf{Hom}}

\newcommand{\eps}{\varepsilon}

% Math Operators

% Special Commands



% DOC
\begin{document}
	\assignment{8}

\begin{problem}[84]
	A diagram of the form
	% https://q.uiver.app/#q=WzAsMyxbMCwwLCJBIl0sWzIsMCwiQiJdLFs0LDAsIkMiXSxbMCwxLCJnIiwyXSxbMSwyLCJlIl0sWzAsMSwiZiIsMCx7Im9mZnNldCI6LTN9XSxbMSwwLCJ0IiwyLHsib2Zmc2V0Ijo0LCJjdXJ2ZSI6Mn1dLFsyLDEsInMiLDIseyJvZmZzZXQiOjIsImN1cnZlIjoyfV1d
	\[\begin{tikzcd}
		A && B && C
		\arrow["g"', from=1-1, to=1-3]
		\arrow["f", shift left=3, from=1-1, to=1-3]
		\arrow["t"', shift right=4, curve={height=12pt}, from=1-3, to=1-1]
		\arrow["e", from=1-3, to=1-5]
		\arrow["s"', shift right=2, curve={height=12pt}, from=1-5, to=1-3]
	\end{tikzcd}\]
	
	such that $ef = eg, es = 1_C, se = gt, ft = 1_B$ is called a split coequalizer. Show that if we have a split coequalizer, that $e$ is the coequalizer of f and g in the colimit sense.
\end{problem}
\begin{solution}
	$e$ is a clearly a cocone over the diagram $\bullet \rightrightarrows \bullet$. We must show that it is a colimit. Let $x: B \to D$ be a morphism such that $xf = xg$ (another cocone). Then, $xs$ is one morphism from $C$ to $D$. We would like this morphism to be our unique factorization. 
	
	First, we need to show that any $x$ factors as $xs e$. However, we know $xse = xgt = xft = x$, and so that's taken care of. Hence, if we can show that $xs$ is unique, we will be done. 
	
	Let $y: C \to D$ be another such factorization ($x = ye$). Then, $y = yes = xs$. So, $e$ is a limit cocone, hence a colimit of the diagram, hence a coequalizer, and we are done! \qed
	
	(The relevant picture)
	% https://q.uiver.app/#q=WzAsNCxbMCwwLCJBIl0sWzIsMCwiQiJdLFs0LDAsIkMiXSxbNCwxLCJEIl0sWzAsMSwiZyIsMl0sWzEsMiwiZSJdLFswLDEsImYiLDAseyJvZmZzZXQiOi0zfV0sWzEsMCwidCIsMix7Im9mZnNldCI6NCwiY3VydmUiOjJ9XSxbMiwxLCJzIiwyLHsib2Zmc2V0IjoyLCJjdXJ2ZSI6Mn1dLFsxLDMsIngiLDFdLFsyLDMsIlxcZXhpc3RzISB4cyIsMCx7InN0eWxlIjp7ImJvZHkiOnsibmFtZSI6ImRhc2hlZCJ9fX1dXQ==
	\[\begin{tikzcd}
		A && B && C \\
		&&&& D
		\arrow["g"', from=1-1, to=1-3]
		\arrow["f", shift left=3, from=1-1, to=1-3]
		\arrow["t"', shift right=4, curve={height=12pt}, from=1-3, to=1-1]
		\arrow["e", from=1-3, to=1-5]
		\arrow["x"{description}, from=1-3, to=2-5]
		\arrow["s"', shift right=2, curve={height=12pt}, from=1-5, to=1-3]
		\arrow["{\exists! xs}", dashed, from=1-5, to=2-5]
	\end{tikzcd}\]


	
\end{solution}


\begin{problem}[86]
	Let $f, g: A \to B$. Let $P$ be the pullback of $f$ and $g$, and let $E$ be the equalizer of $f$ and $g$. Must $P \iso E$?
\end{problem}
\begin{solution}
% https://q.uiver.app/#q=WzAsNyxbNSwzLCJBIl0sWzMsNSwiQSJdLFs1LDUsIkIiXSxbMywzLCJQIl0sWzIsMiwiRSJdLFswLDIsIlAiXSxbMiwwLCJQIl0sWzAsMiwiZiIsMV0sWzEsMiwiZyIsMV0sWzMsMSwiXFxwaV9nIiwxXSxbMywwLCJcXHBpX2YiLDFdLFszLDIsIiIsMSx7InN0eWxlIjp7Im5hbWUiOiJjb3JuZXIifX1dLFs0LDAsImUiLDFdLFs0LDEsImUiLDFdLFs0LDMsIlxcZXhpc3RzISBlX1AiLDEseyJzdHlsZSI6eyJib2R5Ijp7Im5hbWUiOiJkYXNoZWQifX19XSxbNSwxLCJcXHBpX2ciLDFdLFs1LDQsIlxcZXhpc3RzISBwX2ciLDEseyJzdHlsZSI6eyJib2R5Ijp7Im5hbWUiOiJkYXNoZWQifX19XSxbNiw0LCJcXGV4aXN0cyEgcF9mIiwxLHsic3R5bGUiOnsiYm9keSI6eyJuYW1lIjoiZGFzaGVkIn19fV0sWzYsMCwiXFxwaV9mIiwxXV0=
\[\begin{tikzcd}
	&& P &&& \\
	\\
	P && E \\
	&&& P && A \\
	\\
	&&& A && B
	\arrow["{\exists! p_f}"{description}, dashed, from=1-3, to=3-3]
	\arrow["{\pi_f}"{description}, from=1-3, to=4-6]
	\arrow["{\exists! p_g}"{description}, dashed, from=3-1, to=3-3]
	\arrow["{\pi_g}"{description}, from=3-1, to=6-4]
	\arrow["{\exists! e_P}"{description}, dashed, from=3-3, to=4-4]
	\arrow["e"{description}, from=3-3, to=4-6]
	\arrow["e"{description}, from=3-3, to=6-4]
	\arrow["{\pi_f}"{description}, from=4-4, to=4-6]
	\arrow["{\pi_g}"{description}, from=4-4, to=6-4]
	\arrow["\lrcorner"{anchor=center, pos=0.125}, draw=none, from=4-4, to=6-6]
	\arrow["f"{description}, from=4-6, to=6-6]
	\arrow["g"{description}, from=6-4, to=6-6]
\end{tikzcd}\]

This is the relevant image to keep in mind. If, in particular, $\pi_g = \pi_f$, then we trivially get an isomorphism between $P$ and $E$, by uniqueness of morphisms and choosing the map $1_P$. However, this is not always the case.

Take, for instance, $B = A = \{1, 2\}, f(1) = 1, f(2) = 2, g(1) = 1, g(2) = 1$ in $\Set$. Then $P = \{(1, 1), (1, 2)\}$ (with the standard projection maps) but $E = \{1\}$. Then, $P \ncong E$, as the two sets have different cardinalities. \qed


\end{solution}


\end{document}
